\section{Introduction}

Anonymous credential schemes let a user prove they hold certain attributes
without revealing who they are or anything about their undisclosed attribute
values. The two main security properties are unlinkability, so two presentations
of the same user cannot be tied to one another, and unforgeability, so a user
can only successfully attest to having attribute values that they truly
possess~\cite{EC:CamLys01}.

Early anonymous credential schemes assume a single issuer certifying all
attributes, creating a single entity responsible for all issuance, and able to
learn all attribute values of every user. Later schemes distribute issuance
across multiple authorities or delegate it to sub-issuers, but the issuer,
though now decentralised, remains conceptually responsible for all attributes.
This reduces trust in any one party, but does not reflect real-world scenarios,
where distinct entities independently certify one or a small subset of
attributes.

In decentralised identity systems, different authorities each typically certify
a small number of attributes: age, residency, citizenship, possession of a
driving licence, possession of a professional qualification. Credential systems
with distributed issuance still require the issuers to coordinate. With many
issuers worldwide, each relevant to only some users, that coordination is
impractical, and between issuers with no relationship to one another it may be
impossible.

\paragraph{Contributions.} We give an aggregate anonymous credential scheme for this setting: credentials
from independent, uncoordinated issuers combine into one compact credential,
provable in one proof.
\begin{outline}
  \item the key observation behind the scheme, in one sentence: tying credentials from unrelated issuers to one user is already half done in every scheme, by whoever certifies users' keys
  \item leads into the next paragraph (the registration authority)
\end{outline}

A registration authority handles registration only, binding users' keys to
fresh tags. Every credential a user later obtains carries their tag. To present
attributes, users aggregate the relevant credentials and prove, in zero
knowledge, that they know the key, hold a valid aggregate credential for the
disclosed values, and that both share a tag. The tag-based binding is what stops
different users' credentials being pooled, and lets honest users aggregate
alone, without contacting any issuer again.  The registration authority is not a
new party that we add to the model. Every anonymous credential scheme already
assumes that users' keys are certified, so that credentials cannot be issued to
keys that users do not control. That assumption is usually left implicit, or
discharged by the issuer itself. We instead name the role, give it exactly one
job, and give it no other capability. It binds a key to a tag at registration
and does nothing else. It never sees an attribute value, is never contacted
again by the user or by any issuer, and holds no state beyond the tags it has
assigned. Making the role explicit is what makes the scheme efficient. Because
every credential carries the same tag, a user can aggregate credentials from
issuers that have never communicated, and a verifier can check the result
without any proof that the credentials belong to the same person. The
alternative, in a model with no registration authority, is to prove that binding
during presentation, which is exactly the cost we avoid. The gain here comes
from describing a participant that the setting already contains, rather than
from removing one.

\subsection*{Related work}

\paragraph{Single-issuer anonymous credential systems.} An early anonymous
credential system issues signatures on committed values, with security under the
strong RSA (sRSA) assumption~\cite{EC:CamLys01,SCN:CamLys02}. A later pairing-based
construction works under the LRSW assumption, though signature size and
verification complexity grow linearly with the number of attribute values
signed~\cite{C:CamLys04}. An extension supports a constrained form of
aggregation, but only in settings where all values are known ahead of
time~\cite{FC:LeeLeeYun13}.

Credentials built on the BBS+ signature scheme support selective disclosure,
where a user reveals some of their attribute values and hides the rest, in a
single-issuer setting, with verification linear in the number of revealed
attributes, and form the basis of the ongoing standardisation effort of the
Internet Engineering Task Force~\cite{EC:BonBoy04a,SCN:AuSusMu06}.

Further improvements targeted efficiency. The Pointcheval-Sanders signature
scheme introduces randomisable signatures with sequential aggregation, and
proves unforgeability in the GGM~\cite{RSA:PoiSan16}.
These systems operate in a centralised, single-issuer setting where one
authority certifies all attributes. More recently, structure-preserving
signatures on equivalence classes enable very efficient privacy-preserving
proofs, with security proven in the GGM~\cite{JC:FucHanSla19}. Unlinkable
proofs of knowledge over a subset of signed messages, building on the
Pointcheval-Sanders signature scheme, are also proven secure in the
GGM~\cite{PKC:Sanders20}.

\paragraph{Multi-issuer anonymous credential systems.} Other work decentralises
the issuer, delegating issuance to sub-issuers or distributing issuance
responsibilities among several
issuers~\cite{TCC:BCKL08,CCS:CamDriDub17,NDSS:SABMD19}. In these cases, the now
decentralised issuer is still responsible for all attributes. While this reduces
the trust placed in a central authority, it does not fully reflect real-world
scenarios where different authorities independently certify distinct attributes.

A multi-issuer credential system supports selective disclosure and
pseudonyms~\cite{AC:CDHK15}. It is proven secure in the universal composability
framework~\cite{FOCS:Canetti01}. Its construction is built from
several distinct building blocks, each analysed under its own assumptions,
including $n$-BSDH and $n$-RootDH.

\paragraph{Aggregate anonymous credentials.} Aggregate credentials aim to
combine independently issued credentials into a single compact credential, while
maintaining user privacy during presentations. Indexed aggregate signature
schemes support anonymous credentials this way, with unforgeability under the
$q$-Asymmetric Double Hidden Strong Diffie-Hellman ($q$-ADHSDH) assumption, and
anonymity under the decisional Diffie-Hellman (DDH) assumption in the random
oracle model~\cite{canard2011anonymous}.

\begin{outline}
  \item H\'ebant and Pointcheval: ART-Sign, aggregate signatures with randomisable tags, and multi-authority credentials from it, constant showing cost, traceability by a designated authority (cite SCN:HebPoi22); closest scheme to ours
  \item their tag: square Diffie-Hellman tuple $(h, h^x, h^{x^2})$, $h$ derived from the user's identity, $x$ generated jointly with a certification authority; the tag is also the user's public key
  \item signatures are linearly homomorphic: two signatures by one issuer under one tag and index let anyone sign any value for that index, so each issuer signs at most once per tag and index and keeps a per-user record of signed indices
  \item bounded variant: replaces the record with a limit on messages per tag, which the issuer still counts, and an issuer key that grows with the limit
\end{outline}

\begin{outline}
  \item their certification authority plays the same part as our registration authority: both are in the certified-key setting, both bind each user to one tag at registration
  \item the difference is what the tag is: in ART-Sign the tag is the public key and its validity is proven again at every showing; ours is a separate value the registration authority signs together with the key
  \item consequences for us: tag has no structure of its own, issuers keep no record, presentations prove knowledge of the registration signature instead
  \item so in the certified-key setting their scheme is close to ours, and separating tag from key is what lets our issuers keep no state
  \item ART-Sign offers traceability, ours does not
\end{outline}

\paragraph{Other multi-authority credentials.}
\begin{outline}
  \item Mir et al.: aggregate signatures with versatile randomisation, multi-authority credentials that also hide which issuers certified the attributes (cite CCS:MBGLS23)
  \item issuer hiding for single credentials (cite CANS:BEKRS21, EPRINT:SanTra23)
  \item credentials from existing identity documents with zk-SNARKs (cite SP:RWGM23)
  \item anonymous authentication decentralised across authorities for global identities (cite EPRINT:Anada19)
  \item a single issuer distributed among several parties without trusted setup (cite EPRINT:GGPR25)
\end{outline}

\begin{outline}
  \item short glossary of the assumptions named above: $q$-ADHSDH, $n$-BSDH, $n$-RootDH are $q$-type (statement grows with the number of signing queries)
  \item strong RSA: hard to find any root of a given value modulo an RSA modulus
  \item LRSW: interactive assumption in prime-order groups (cite SAC:LRSW99)
\end{outline}

\begin{outline}
  \item our scheme: same asymptotic costs as ART-Sign
  \item issuers are stateless and need no bound on the number of attributes in advance; tags are bound to keys by the registration authority, not the keys themselves
  \item both schemes are proven in the GGM, so the comparison is about issuer state and how tags are bound to users (not about assumptions)
  \item point to \Cref{tab:agg-creds-related} for the comparison with single- and multi-issuer schemes
\end{outline}

\begin{outline}
  \item issuer state, said in prose now that the table has no column for it: ART-Sign issuers must keep per-user state, a record of the indices they have signed under each tag, since two signatures on the same tag and index let anyone sign any value for that index (cite SCN:HebPoi22)
  \item our issuers keep none: they check the registration signature and sign, and remember nothing about the user
  \item CL02, CL04 and BBS+ are single-issuer, so the question does not arise there; CDHK15 issuers keep no per-user record either (cite AC:CDHK15); CL11 does not consider independent issuers (cite canard2011anonymous)
\end{outline}

\begin{table}[htbp] \centering
  \begin{tabular}{l|c|c|c|c}
    \toprule Scheme & Assumptions & Multi-issuer & $|$cred$|$ & Prover cost \\ \midrule
    CL02~\cite{SCN:CamLys02} & sRSA & No & $O(1)$ & $O(n)$ \\
    CL04~\cite{C:CamLys04} & LRSW & No & $O(n)$ & $O(n)$ \\
    BBS+ (BB04, ASM06)~\cite{EC:BonBoy04a,SCN:AuSusMu06} & $q$-SDH & No & $O(1)$ & $O(n)$ \\
    CL11~\cite{canard2011anonymous} & $q$-ADHSDH & Partial & $O(n)$ & $O(k)$ \\
    CDHK15~\cite{AC:CDHK15} & $n$-BSDH, $n$-RootDH & Yes & $O(1)$ & $O(k)$ \\
    HP22~\cite{SCN:HebPoi22} & GGM & Yes & $O(1)$ & $O(k)$ \\
    Our scheme & GGM & Yes & $O(1)$ & $O(k)$ \\
\bottomrule
  \end{tabular}
  \caption[Comparison with existing anonymous credential schemes]{Comparison with
existing single and multi-issuer anonymous credential schemes. $n$ is the total
number of attributes a scheme supports, $k$ the number of attributes disclosed in
a presentation, and $|$cred$|$ the size of a full aggregated credential. The
assumptions of the other schemes are those of the cited papers. Canard and
Lescuyer's scheme aggregates signatures under a shared index, but does not
consider independent issuers in the sense we
require~\cite{canard2011anonymous}.} \label{tab:agg-creds-related} \end{table}

In contrast to classical anonymous credential
schemes, we explicitly consider a decentralised setting in which each attribute
is certified by a distinct issuer, and credentials issued by different
authorities are combined into one single compact credential.  Aggregation is
performed locally by the user without interaction between issuers. The scheme
supports selective disclosure of attribute values.

\paragraph{Layout.} The rest of the chapter is organised as follows. \Cref{sec:agg-creds-definitions} states the problem and
formalises it as an ideal functionality.  \Cref{sec:agg-creds-construction}
gives an abstract construction from generic primitives.
\Cref{sec:agg-creds-instantiation} gives a concrete instantiation, including our
tag-based aggregate signature scheme, which \Cref{sec:agg-creds-security} proves
secure.  \Cref{sec:agg-creds-evaluation} evaluates performance, and
\Cref{sec:agg-creds-summary} summarises what the chapter shows and where it leads.
